\section{Sincronização de posição e atitude}
\label{sec:sincronizacao}
A sincronização das posições e das atitudes de cada espaçonave corresponde à condição em que o veículo \emph{chaser}, além de atingir a posição desejada em relação ao \emph{target}, também apresenta orientação coincidente com a da espaçonave alvo.
Essa exigência permite que o alinhamento geométrico entre os sistemas seja preservado, que é uma condição para as operações de \emph{rendezvous and docking/berthing} ocorra com sucesso.
Sendo assim, dois critérios principais foram adotados:

(i) O erro angular relativo entre chaser e target deve ser inferior a $1^\circ$, assegurando que a orientação seja suficientemente precisa para evitar desalinhamentos críticos,
e (ii) a posição do chaser deve permanecer dentro do corredor seguro de aproximação, para que possa reduzir os riscos de colisão e permitir manobras corretivas em caso de falha parcial do controle.

A sincronização de posição e atitude está diretamente relacionada à segurança da aproximação, pois garante que a trajetória de entrada do chaser seja compatível com os limites estabelecidos para o acoplamento.
Além disso, está vinculada à operação de \emph{berthing}, pois o alinhamento entre os eixos estruturais é indispensável para o sucesso do engate mecânico entre a ferramenta e a porta de \emph{docking}.

Por fim, a avaliação do desempenho do sistema de sincronização é conduzida a partir de cenários simulados, que permitem verificar a robustez da lei de controle frente a perturbações paramétricas.
Essa conexão entre requisitos teóricos e simulações práticas fornece a base para validar os critérios de projeto que foram estabelecidos.